\chapter{CALCULOS JUSTIFICATIVOS}

\section{Alcance del calculo}

El calculo se desarrolla para una vivienda unifamiliar de tres niveles con suministro monofasico de 220 V, 60 Hz. La propuesta se basa en los layouts arquitectonicos canonicos y en el modelo electrico vigente representado en las laminas IE-02, IE-03 e IE-04.

La sectorizacion se organiza en 7 circuitos independientes: alumbrado y tomacorrientes por nivel, mas un circuito independiente de tomacorrientes especiales de cocina en el primer piso. No se consideran motores, bombas ni cargas especiales ajenas al levantamiento real de puntos.

\section{Premisas de calculo}

\begin{itemize}
  \item Tension nominal: 220 V monofasico.
  \item Factor de potencia (cos $\phi$) adoptado: 0.90.
  \item Potencia de diseno para circuitos de alumbrado: 500 W por nivel, adoptada como reserva academica por encima de la carga LED instalada.
  \item Potencia de diseno para circuitos de tomacorrientes: 1000 W a 1500 W por circuito, segun uso y concentracion de puntos.
  \item Factor de demanda para circuitos de alumbrado: 1.00.
  \item Factor de demanda para tomacorrientes generales en segundo y tercer piso: 0.70.
  \item Factor de demanda para tomacorrientes del primer piso y cocina: 1.00, por criterio conservador.
  \item Proteccion diferencial de 30 mA para todos los circuitos de tomacorrientes.
\end{itemize}

\section{Formulas empleadas}

\begin{equation}
P_{inst} = \sum P_i
\end{equation}

\begin{equation}
P_{dem} = \sum (P_i \times f_{d,i})
\end{equation}

\begin{equation}
I_{dem} = \frac{P_{dem}}{V \times fp}
\end{equation}

Donde $P_{inst}$ es la potencia instalada, $P_{dem}$ es la potencia demandada, $f_{d,i}$ es el factor de demanda del circuito $i$, $V$ es la tension nominal (220 V) y $fp$ es el factor de potencia (0.90).

\section{Cuadro de cargas por circuito}

\begin{table}[H]
\centering
\footnotesize
\setlength{\tabcolsep}{2.5pt}
\caption{Cuadro de cargas por circuito}
\label{tab:cuadro-cargas}
\begin{tabularx}{\textwidth}{c L{3.3cm} r c L{3.1cm} c L{2.2cm}}
\toprule
\textbf{Cto.} & \textbf{Descripcion} & \textbf{Pot. (W)} & \textbf{F.D.} & \textbf{Conductor} & \textbf{ITM} & \textbf{ID} \\
\midrule
C1 & Alumbrado primer piso & 500 & 1.00 & 3 x 1.5 mm2 Cu - PVC 3/4" & 2P-10A & ID general \\
C2 & Tomacorrientes generales primer piso & 1000 & 1.00 & 3 x 2.5 mm2 Cu - PVC 3/4" & 2P-16A & 2P-25A/30mA \\
C3 & Tomacorrientes especiales de cocina primer piso & 1500 & 1.00 & 3 x 2.5 mm2 Cu - PVC 3/4" & 2P-20A & 2P-25A/30mA \\
C4 & Alumbrado segundo piso & 500 & 1.00 & 3 x 1.5 mm2 Cu - PVC 3/4" & 2P-10A & ID general \\
C5 & Tomacorrientes generales segundo piso & 1500 & 0.70 & 3 x 2.5 mm2 Cu - PVC 3/4" & 2P-16A & 2P-25A/30mA \\
C6 & Alumbrado tercer piso & 500 & 1.00 & 3 x 1.5 mm2 Cu - PVC 3/4" & 2P-10A & ID general \\
C7 & Tomacorrientes generales tercer piso & 1500 & 0.70 & 3 x 2.5 mm2 Cu - PVC 3/4" & 2P-16A & 2P-25A/30mA \\
\midrule
\textbf{TOTAL} & \textbf{Demanda consolidada} & \textbf{7000} & -- & Acometida: 2 x 10 mm2 Cu + PE & Gral: 2P-40A & ID: 2P-40A/30mA \\
\bottomrule
\end{tabularx}
\end{table}

\section{Levantamiento de cargas por circuito}

\begin{table}[H]
\centering
\small
\setlength{\tabcolsep}{3pt}
\caption{Levantamiento de cargas por circuito}
\label{tab:levantamiento-cargas}
\begin{tabularx}{\textwidth}{c Y L{3.2cm} r c}
\toprule
\textbf{Item} & \textbf{Circuito / carga} & \textbf{Puntos vinculados} & \textbf{Pot. Inst. (W)} & \textbf{F. Demanda} \\
\midrule
1 & C1 - Alumbrado primer piso & 7 luminarias / 5 interruptores & 500 & 1.00 \\
2 & C2 - Tomacorrientes generales primer piso & 4 tomacorrientes dobles & 1000 & 1.00 \\
3 & C3 - Tomacorrientes especiales de cocina primer piso & 4 tomacorrientes dobles & 1500 & 1.00 \\
4 & C4 - Alumbrado segundo piso & 6 luminarias / 5 interruptores & 500 & 1.00 \\
5 & C5 - Tomacorrientes generales segundo piso & 8 tomacorrientes, incluido 1 protegido & 1500 & 0.70 \\
6 & C6 - Alumbrado tercer piso & 6 luminarias / 6 interruptores & 500 & 1.00 \\
7 & C7 - Tomacorrientes generales tercer piso & 8 tomacorrientes, incluido 1 protegido & 1500 & 0.70 \\
\bottomrule
\end{tabularx}
\end{table}

\section{Interpretacion de la demanda maxima}

La potencia instalada total es 7000 W. Aplicando los factores de demanda del modelo electrico canonico se obtiene una demanda maxima estimada de 6100 W. Con tension de 220 V y factor de potencia 0.90:

\begin{equation}
I_{dem} = \frac{6100}{220 \times 0.90} = 30.81\text{ A}
\end{equation}

Con este valor de corriente de demanda se dimensiona el alimentador general con conductor de cobre de 10 mm2 y un interruptor termomagnetico principal de 2P-40A.

\section{Tablero general y tableros de distribucion}

\begin{table}[H]
\centering
\small
\caption{Propuesta de tableros}
\label{tab:tableros-distribucion}
\begin{tabularx}{\textwidth}{L{1.8cm}L{4.0cm}L{4.2cm}Y}
\toprule
\textbf{Elemento} & \textbf{Ubicacion propuesta} & \textbf{Circuitos asociados} & \textbf{Criterio} \\
\midrule
TG-01 & Primer piso, pasadizo longitudinal & C1, C2 y C3 & Tablero principal accesible y rotulado \\
TD-02 & Segundo piso, zona central/hall & C4 y C5 & Subtablero para el control del 2do nivel \\
TD-03 & Tercer piso, pasadizo longitudinal & C6 y C7 & Subtablero para el control del 3er nivel \\
\bottomrule
\end{tabularx}
\end{table}

\section{Protecciones preliminares}

\begin{table}[H]
\centering
\small
\setlength{\tabcolsep}{3pt}
\caption{Protecciones preliminares por circuito}
\label{tab:protecciones-preliminares}
\begin{tabularx}{\textwidth}{l Y Y Y Y}
\toprule
\textbf{Cto.} & \textbf{Uso} & \textbf{Interruptor} & \textbf{Conductor} & \textbf{Observacion} \\
\midrule
General & Alimentacion principal & 2P-40A & 3 x 10 mm2 Cu + PE & Llave general termomagnetica \\
- & ID General & 2P-40A / 30mA & - & Interruptor diferencial general TG \\
C1 & Alumbrado primer piso & 2P-10A & 3 x 1.5 mm2 Cu & Llave termomagnetica \\
C2 & Tomacorrientes generales 1er piso & 2P-16A / ID-25A & 3 x 2.5 mm2 Cu & Con protector diferencial \\
C3 & Tomacorrientes de cocina 1er piso & 2P-20A / ID-25A & 3 x 2.5 mm2 Cu & Con protector diferencial \\
C4 & Alumbrado segundo piso & 2P-10A & 3 x 1.5 mm2 Cu & Llave termomagnetica \\
C5 & Tomacorrientes del segundo piso & 2P-16A / ID-25A & 3 x 2.5 mm2 Cu & Con protector diferencial \\
C6 & Alumbrado tercer piso & 2P-10A & 3 x 1.5 mm2 Cu & Llave termomagnetica \\
C7 & Tomacorrientes del tercer piso & 2P-16A / ID-25A & 3 x 2.5 mm2 Cu & Con protector diferencial \\
\bottomrule
\end{tabularx}
\end{table}

\section{Coordinacion conductor-proteccion}

La seleccion preliminar mantiene el criterio basico:

\begin{equation}
I_b \leq I_n \leq I_z
\end{equation}

Donde $I_b$ es la corriente de diseno, $I_n$ es la corriente nominal del interruptor e $I_z$ es la capacidad admisible del conductor. Con calibres de 1.5 mm2 ($I_z \approx 15\text{ A}$) y de 2.5 mm2 ($I_z \approx 20\text{ A}$), las llaves de 10A, 16A y 20A garantizan la coordinacion termica de los circuitos derivados. El alimentador de 10 mm2 mantiene holgura para la demanda de 30.81 A con proteccion general de 40 A.

\section{Sistema de puesta a tierra}

La propuesta incluye conductor de proteccion de tierra en todos los tomacorrientes y barra de conexion a tierra en cada tablero. La conexion final se realiza a un pozo de tierra vertical con electrodo de cobre de 5/8" x 2.40 m y aditivo de conductividad gel, disenado para una resistencia de puesta a tierra menor a 15 Ohms.

\section{Resumen de puntos instalados}

\begin{table}[H]
\centering
\small
\caption{Resumen de puntos electricos por tipo}
\label{tab:resumen-puntos}
\begin{tabularx}{\textwidth}{L{3.5cm} c c c c}
\toprule
\textbf{Nivel} & \textbf{Luminarias} & \textbf{Interruptores} & \textbf{Tomacorrientes} & \textbf{Total} \\
\midrule
Primer piso & 7 & 5 & 8 & 20 \\
Segundo piso & 6 & 5 & 8 & 19 \\
Tercer piso & 6 & 6 & 8 & 20 \\
\midrule
\textbf{Total} & \textbf{19} & \textbf{16} & \textbf{24} & \textbf{59} \\
\bottomrule
\end{tabularx}
\end{table}

\section{Observaciones de auditoria}

\begin{itemize}
  \item \textbf{Cargas especiales eliminadas:} no se incluyen electrobombas ni motores especiales que no correspondan a la solicitud final del diseno de puntos.
  \item \textbf{Circuitos resultantes:} 7 circuitos generales (C1 a C7) que cubren alumbrado y tomacorrientes por nivel, mas la cocina independiente.
  \item \textbf{Puntos instalados:} 19 luminarias, 16 interruptores y 24 tomacorrientes, para un total de 59 puntos fisicos vinculados a los 7 circuitos derivados.
\end{itemize}
